\chap{绪论：智能软件开发的问题与愿景}
\label{sec:intro}
\subsection{从完成一次任务到长期发展一个系统}
软件开发中的一个局部任务往往容易描述：增加一个按钮、修复一个接口，或让某个测试通过。一个长期存在的软件系统却同时承载许多尚未完成的事情。已有用户依赖它的行为，新的需求持续进入，多个模块共同完成同一任务，过去的设计理由未必能从代码直接读出。开发者必须在变化中决定哪些工作值得做、怎样分配、如何确认结果，以及何时让用户实际用到成果。

当 AI 可以阅读代码、调用工具和修改实现时，这些工程判断可以逐步进入它的工作范围。由此形成的开发方式，需要同时考虑个体 Agent 的能力与整个工作系统的组织。一次成功修改证明了一个具体能力；持续接收目标、维持质量、吸收反馈并改善后续工作的能力，则需要在更长时间内观察。本文把这一持续过程作为讨论对象。

这里的“智能软件开发”指 AI 深度参与并自主承担工程判断的开发方式。所开发的软件可以是普通业务系统，也可以包含模型服务、智能助手或自动化模块。产物中是否使用 AI，与开发过程是否具有智能组织能力，是两个可以分别讨论的维度。本文重点讨论后者，并在最后讨论二者进一步结合的可能。

一个基本场景可以说明其要求。使用者希望任务平台支持可靠取消长时间运行的工作。实现涉及前端反馈、任务状态、调度、执行器和资源释放。只增加一个“取消”按钮无法完成目标；只向用户解释若干接口也不足以确认体验。开发系统需要理解取消的含义，找到责任模块，处理并发条件，验证用户看到的结果，再把这些结果以便于人判断的方式呈现。后续用户可能进一步要求保留部分输出，原有设计还要继续演化。

\subsection{三个相互作用的目的}
本文把智能开发系统的目的概括为三项：让人的意图更准确地进入 AI 的工作，让 AI 更自主地完成所委派的工作，以及让 AI 的成果与机制更容易被人理解和控制。三项目的相互影响。难以理解的汇报会降低反馈质量；含糊的意图会增加实现返工；执行缺乏连续性又会迫使人反复解释背景。

因此，减少人类控制成本并不等同于减少消息数量。一次短消息可能包含需要长时间调查的判断，一张复杂仪表盘也可能把大量解释工作转移给人。更合适的问题是：为了让系统继续朝期望方向发展，人需要花多少时间建立理解、比较选项、补充信息和纠正偏差？系统能否把已经表达过的偏好用于后续判断，同时在真正重要的分歧上提供有用的参与入口？

自主性也应联系结果理解。Agent 能连续调用许多工具，说明它具备执行能力；开发系统还需要判断这些动作是否服务于当前目标，是否产生可用成果，以及失败之后是否完成修复。把任务投递给另一个成员只是委派的开始。对人有意义的工作单元通常延续到成果可体验、影响可解释、遗留问题可定位的阶段。

\fig{01-overview}{智能开发系统的整体关系。人的意图进入语义约束与自主组织；开发、验证和使用产生证据；交互投影将成果与操作带回人的工作情境。实线表示主要工作依赖，虚线表示反馈。}{overview}

\subsection{实践动机与研究范围}
Portal 是一个围绕项目、知识、持续 Agent 和协作过程构建的系统。其开发实践提供了本文的问题来源。项目曾尝试把大量设计点、改进点和实现差异集中呈现在 WALL 界面中，使人可以监督系统的发展。随着点的数量增加，逐项查看的负担显著上升。更深的问题是，单条简短文字往往无法提供足够背景，人仍需请求解释才能判断它的含义。

另一个实践观察是，直接使用软件往往更容易产生明确反馈。操作界面同时提供目标、对象、行为和结果，人能够指出具体哪里不符合预期。此经验启发了一个设计方向：为架构、机制和未来方案也构造可探索的表达，让人能在相同情境中理解并采取行动。本文将这一方向与团队组织、持续交付和验证联系起来。

这些观察来自一个持续变化的项目及其使用者经验，不能据此推断所有团队都具有相同问题强度。本文采用概念性设计研究的写法：从明确的实践动机出发，提出机制及其作用解释，再给出可检验的评价方案。文中的过程图和界面图均为作者绘制的设计示意；第\ref{sec:case}节的完整工作过程是构造案例，用于检验框架内部的一致性。

\subsection{与既有研究的联系}
软件模块化研究强调将设计决定与职责组织起来，使局部变化具有可控制的影响范围\cite{parnas1972}。组织沟通与系统结构之间的联系也已有长期讨论\cite{conway1968}。这些思想对 Agent 团队仍有意义，但持续身份、可复制上下文和低成本临时成员提供了新的组织手段。本文关注这些手段怎样服务于长期维护，以及怎样避免协调开销抵消并行收益。

语言模型 Agent 研究展示了工具交互、环境反馈和多成员协作的多种实现方向。ReAct 将推理与行动交替组织\cite{react}；SWE-agent 研究面向 Agent 的计算机接口如何支持软件工程任务\cite{sweagent}；MetaGPT 和 AutoGen 分别探索协作工作流与多 Agent 会话组织\cite{metagpt,autogen}。这些工作为自主执行提供了参考。本文进一步讨论任务完成之后的持续演化、人的方向控制，以及模型升级后组织方式的适应。

在人机交互方面，逐层查看、按需获取细节与混合主动性交互已有研究基础\cite{shneiderman1996,horvitz1999}。人机 AI 交互指南也关注能力说明、纠错和随时间变化的行为\cite{amershi2019}。本文把这些原则应用于开发系统：由 Agent 结合当前工作生成表达，允许人在对应对象上参与，并将交互反馈持续用于更新软件和表达方式。

\begin{table}[htbp]\centering\small
\caption{相关研究与本文问题的联系}\label{tab:related}
\begin{tabularx}{\linewidth}{p{2.4cm}YY}\toprule
研究方向 & 代表性研究提供的认识 & 本文进一步讨论的问题\\\midrule
模块化与组织 & 设计决定的分解及沟通结构的影响\cite{parnas1972,conway1968} & 持续 Agent 的职责、上下文和经验怎样沿边界积累\\
自主编码 & 工具交互、仓库操作与任务结果评价\cite{sweagent,swebench} & 多轮需求变化、交付修复与人的持续引导\\
多 Agent 协作 & 工作流和会话能够组织多成员行动\cite{metagpt,autogen} & 何时委派、如何调整职责、怎样控制协调成本\\
人机交互 & 逐步探索、混合主动性及纠错原则\cite{shneiderman1996,horvitz1999,amershi2019} & 如何根据当前开发工作生成并维护理解与控制界面\\\bottomrule
\end{tabularx}\end{table}

这些研究分别提供了重要基础。本文将它们联系到同一个长期工作对象：一个需要持续维护并接受人类引导的软件系统。其提出的机制并非全部为新发明，研究价值主要在于明确组合关系、成本条件和可评价的效果。例如，较好的模块边界可以减少交付依赖，也能让架构投影更容易理解；更准确的人类反馈则可能减少下一轮实现返工。

文章也避免将传统软件工程概括为单一固定流程。成熟的人类团队同样使用持续交付、探索式测试和按职责分工。这里的变化是，更多选择这些方法并调整组织的判断可以由 AI 持续承担。需要研究的是在什么条件下这种委派有效，以及系统怎样依据实践继续改进。

\subsection{文章贡献与阅读路线}
本文提出一个连接组织、工程和交互的整体框架，具体贡献包括：定义面向能力放大的目标与成本视角；说明持续成员、架构职责和递归委派的关系；提出局部交付与证据驱动验证的组合方式；给出理解与控制共处的工作投影原则及典型指南；设计包含长期变化、人类负担和资源投入的评价方法。

第\ref{sec:framework}节建立概念与整体关系，第\ref{sec:intent}至\ref{sec:validation}节讨论方向、组织、交付和验证，第\ref{sec:interaction}与\ref{sec:guidelines}节讨论面向人的表达，第\ref{sec:case}节串联全过程。第\ref{sec:evaluation}节把设计主张转化为评价问题，最后讨论开放问题。读者可以先从总体框架与案例建立直觉，再回到具体机制理解条件和取舍。
